\documentclass[11pt,a4paper]{article}

\usepackage[margin=2.2cm]{geometry}
\usepackage{amsmath,amssymb,amsthm}
\usepackage{booktabs}
\usepackage{array}
\usepackage{tabularx}
\usepackage[shortlabels]{enumitem}
\usepackage{xcolor}
\usepackage{tikz}
\usepackage{fancyhdr}
\usepackage{setspace}
\usepackage{microtype}
\usepackage{titlesec}
\usepackage{hyperref}
\usepackage{listings}
\usepackage[skins,breakable]{tcolorbox}

\usetikzlibrary{positioning,arrows.meta,calc,fit,shapes.geometric,matrix,decorations.pathreplacing}

\definecolor{ink}{HTML}{222222}
\definecolor{rulegray}{HTML}{444444}
\definecolor{boxline}{HTML}{666666}
\definecolor{boxfill}{HTML}{F4F4F4}
\definecolor{boxtitle}{HTML}{E8E8E8}
\definecolor{g1}{HTML}{D0D0D0}
\definecolor{g2}{HTML}{9A9A9A}
\definecolor{g3}{HTML}{5A5A5A}

\hypersetup{
  colorlinks=true,
  linkcolor=black,
  urlcolor=black,
  pdftitle={MTL 458 Lecture 18: Condition Variables},
  pdfauthor={Minati De, IIT Delhi}
}

\setstretch{1.12}
\raggedbottom
\setlist[itemize]{leftmargin=1.3em,itemsep=0.28em,topsep=0.35em}
\setlist[enumerate]{leftmargin=1.5em,itemsep=0.35em,topsep=0.4em}

\titleformat{\section}
  {\Large\bfseries\color{ink}}{\thesection.}{0.6em}{}
  [{\color{rulegray}\titlerule[0.5pt]}]
\titleformat{\subsection}
  {\large\bfseries\color{ink}}{\thesubsection}{0.55em}{}
\titleformat{\subsubsection}
  {\normalsize\bfseries\color{ink}}{\thesubsubsection}{0.5em}{}

\titlespacing*{\section}{0pt}{1.35em}{0.7em}
\titlespacing*{\subsection}{0pt}{1.05em}{0.45em}

\pagestyle{fancy}
\fancyhf{}
\renewcommand{\headrulewidth}{0.4pt}
\renewcommand{\footrulewidth}{0.3pt}
\fancyhead[L]{\small\color{ink}MTL 458 \textbullet\ Operating Systems}
\fancyhead[R]{\small\color{ink}Lecture 18: Condition Variables}
\fancyfoot[C]{\small\thepage}
\setlength{\headheight}{14pt}

% ---------------------------------------------------------------------------
% Boxes
% ---------------------------------------------------------------------------
\tcbset{
  commonbox/.style={
    colback=boxfill, colframe=boxline, boxrule=0.5pt,
    left=6pt, right=6pt, top=5pt, bottom=5pt,
    fonttitle=\bfseries, coltitle=black, colbacktitle=boxtitle,
    enhanced, sharp corners=south, breakable
  }
}
\newtcolorbox{keybox}[1][]{commonbox, title={#1}}
\newtcolorbox{quizbox}[1][]{commonbox, title={#1}}
\newtcolorbox{solbox}[1][]{commonbox, title={#1}}
\newtcolorbox{warnbox}[1][]{commonbox, title={#1}}
\newtcolorbox{defbox}[1][]{commonbox, title={#1}}

\lstset{
  language=C,
  basicstyle=\ttfamily\small,
  keywordstyle=\bfseries,
  commentstyle=\itshape\color{g3},
  numbers=left, numberstyle=\scriptsize\color{g2}, numbersep=6pt,
  frame=single, rulecolor=\color{boxline},
  backgroundcolor=\color{boxfill},
  xleftmargin=1.6em, framexleftmargin=1.4em,
  columns=fullflexible, keepspaces=true,
  showstringspaces=false, breaklines=true,
  aboveskip=0.6em, belowskip=0.6em,
  morekeywords={cond_t,mutex_t,pthread_t,pthread_cond_t,pthread_mutex_t}
}
\lstdefinestyle{term}{language={}, numbers=none, basicstyle=\ttfamily\footnotesize,
  xleftmargin=0.4em, framexleftmargin=0.2em}

% Start a new page unless at least #1 of vertical space remains
\makeatletter
\newcommand{\keepspace}[1]{\par\penalty-100\begingroup
  \setlength{\dimen@}{#1}\dimen@ii\pagegoal\advance\dimen@ii-\pagetotal
  \ifdim\dimen@>\dimen@ii\ifdim\dimen@ii>\z@\vfil\fi\break\fi\endgroup}
\makeatother

\newcommand{\Tc}[1]{\(T_{c#1}\)}
\newcommand{\Tp}{\(T_p\)}

% Schedule rows: \sched{parent events}{child events}, each a list of position/label
\newcommand{\sched}[2]{%
\begin{tikzpicture}[x=0.85cm, y=0.6cm, every node/.style={font=\small}]
  \node[anchor=east] at (-0.2,1) {\textbf{Parent:}};
  \node[anchor=east] at (-0.2,0) {\textbf{Child:}};
  \draw[g2] (0,1) -- (11,1); \draw[g2] (0,0) -- (11,0);
  \foreach \x/\l in {#1} { \node[fill=white, inner sep=1.5pt, font=\ttfamily\small] at (\x,1) {\l}; }
  \foreach \x/\l in {#2} { \node[fill=white, inner sep=1.5pt, font=\ttfamily\small] at (\x,0) {\l}; }
  \draw[-{Latex}, g3] (0,-0.7) -- (11,-0.7) node[right, font=\scriptsize] {time};
\end{tikzpicture}}

% Pipe buffer: \pipebuf{filled segments a/b}{start}{end}
\newcommand{\pipebuf}[3]{%
\begin{tikzpicture}[x=3.85cm, y=0.42cm, baseline=(base)]
  \useasboundingbox (-0.06,-1.75) rectangle (1.06,2.55);
  \coordinate (base) at (0,0.3);
  \foreach \a/\b in {#1} { \fill[black!50] (\a,0) rectangle (\b,1); }
  \draw[thick] (0,0) rectangle (1,1);
  \draw[-{Latex[length=1.3mm]}] (#2,1.75) node[above, font=\tiny, inner sep=1pt] {start} -- (#2,1.05);
  \draw[-{Latex[length=1.3mm]}] (#3,-0.75) node[below, font=\tiny, inner sep=1pt] {end} -- (#3,-0.05);
\end{tikzpicture}}

\begin{document}

% ===========================================================================
% Title
% ===========================================================================
\begin{center}
{\color{rulegray}\rule{\textwidth}{1.1pt}}\\[0.45em]
{\LARGE\bfseries Operating Systems}\\[0.25em]
{\Large\bfseries MTL 458 \textperiodcentered\ Lecture 18}\\[0.35em]
{\LARGE\bfseries Condition Variables}\\[0.35em]
{\large Minati De}\\[0.15em]
{\normalsize Department of Mathematics, Indian Institute of Technology Delhi}\\[0.35em]
{\color{rulegray}\rule{\textwidth}{1.1pt}}
\end{center}

\vspace{0.4em}
\begin{keybox}[Questions answered in this lecture]
\begin{itemize}
  \item How can threads enforce \textbf{ordering} across operations?
  \item How can \textbf{\texttt{thread\_join()}} be implemented?
  \item How can \textbf{condition variables} be used to support
        \textbf{producer/consumer} apps?
\end{itemize}
\end{keybox}

\tableofcontents
\vspace{0.6em}
\noindent{\color{rulegray}\rule{\textwidth}{0.4pt}}

% ===========================================================================
\section{Concurrency objectives}
% ===========================================================================

\begin{center}
\renewcommand{\arraystretch}{1.25}
\begin{tabularx}{\textwidth}{@{}lXl@{}}
  \toprule
  \textbf{Objective} & \textbf{Example} & \textbf{Solved with} \\
  \midrule
  \textbf{Mutual exclusion} & A and B don't run at the same time & \emph{locks} \\
  \textbf{Ordering} & B runs after A does something & \emph{condition variables} and \emph{semaphores} \\
  \bottomrule
\end{tabularx}
\end{center}

The previous lecture solved mutual exclusion. This lecture is about
\textbf{ordering}.

% ===========================================================================
\section{Condition variables}
% ===========================================================================

There are many cases where a thread wishes to \textbf{check} whether a
\textbf{condition} is true before continuing its execution.

\textbf{Example:} a parent thread might wish to check whether a child
thread has \emph{completed}. This is often called a \texttt{join()}.

\subsection{A parent waiting for its child}

\begin{lstlisting}
void *child(void *arg) {
    printf("child\n");
    // XXX how to indicate we are done?
    return NULL;
}

int main(int argc, char *argv[]) {
    printf("parent: begin\n");
    pthread_t c;
    Pthread_create(&c, NULL, child, NULL); // create child
    // XXX how to wait for child?
    printf("parent: end\n");
    return 0;
}
\end{lstlisting}

What we would like to see here is:
\begin{lstlisting}[style=term]
parent: begin
child
parent: end
\end{lstlisting}

Without waiting, \texttt{parent: end} could be printed before
\texttt{child}.

\keepspace{22\baselineskip}
\subsection{Spin-based approach}

\begin{lstlisting}
int done = 0;

void *child(void *arg) {
    printf("child\n");
    done = 1;
    return NULL;
}

int main(int argc, char *argv[]) {
    printf("parent: begin\n");
    pthread_t c;
    Pthread_create(&c, NULL, child, NULL); // create child
    while (done == 0)
        ; // spin
    printf("parent: end\n");
    return 0;
}
\end{lstlisting}

The child sets a shared flag \texttt{done}; the parent loops until it sees
it. This prints the right output, but:

\begin{warnbox}[Spinning]
This is hugely \textbf{inefficient}: the parent spins and \textbf{wastes CPU
time} (the same problem as spinlocks in the last lecture).
\end{warnbox}

\subsection{How to wait for a condition}

\begin{defbox}[Condition variable]
\begin{itemize}
  \item \textbf{Waiting} on the condition: a condition variable is an
        \textbf{explicit queue} that threads can put themselves on when
        some state of execution is not as desired.
  \item \textbf{Signaling} on the condition: some other thread, \emph{when
        it changes said state}, can wake one of those waiting threads and
        allow them to continue.
\end{itemize}
\end{defbox}

\keepspace{20\baselineskip}
\subsection{Definition and routines}

Declare a condition variable (proper initialization is required):
\begin{lstlisting}[numbers=none]
pthread_cond_t c;
\end{lstlisting}

The POSIX calls:
\begin{lstlisting}[numbers=none]
pthread_cond_wait(pthread_cond_t *c, pthread_mutex_t *m); // wait()
pthread_cond_signal(pthread_cond_t *c);                   // signal()
\end{lstlisting}

\begin{keybox}[What \texttt{wait()} does with the mutex]
\begin{itemize}
  \item The \texttt{wait()} call takes a \textbf{mutex} as a parameter (the
        caller must hold it).
  \item \texttt{wait()} \textbf{releases the lock} and puts the calling
        thread to sleep, as one atomic step. This is what prevents a race
        while a thread is putting itself to sleep: no other thread can
        change the state and signal between the caller's check and its
        sleep (Section~3 shows what goes wrong without this).
  \item When the thread wakes up, it must \textbf{re-acquire the lock}
        before \texttt{wait()} returns.
\end{itemize}
\texttt{signal()} wakes one thread waiting on \texttt{c}; if nobody is
waiting, it does nothing.
\end{keybox}

% ===========================================================================
\section{Implementing join with a condition variable}
% ===========================================================================

In the slides \texttt{Mutex\_lock}, \texttt{Cond\_wait}, etc.\ are wrappers
for the \texttt{pthread\_} calls. Letters in comments label each line so we
can write schedules.

\subsection{Attempt 1}

\begin{center}
\begin{minipage}[t]{0.48\textwidth}
\textbf{Parent:}
\begin{lstlisting}[numbers=none]
void thread_join() {
  Mutex_lock(&m);     // x
  Cond_wait(&c, &m);  // y
  Mutex_unlock(&m);   // z
}
\end{lstlisting}
\end{minipage}\hfill
\begin{minipage}[t]{0.48\textwidth}
\textbf{Child:}
\begin{lstlisting}[numbers=none]
void thread_exit() {
  Mutex_lock(&m);     // a
  Cond_signal(&c);    // b
  Mutex_unlock(&m);   // c
}
\end{lstlisting}
\end{minipage}
\end{center}

\textbf{Example schedule:}
\begin{center}
\sched{1/x, 2.5/y, 9.5/z}{5/a, 6.5/b, 8/c}
\end{center}
\textbf{Works!} The parent waits at \texttt{y} (releasing \texttt{m}); the
child signals at \texttt{b}; the parent wakes, re-acquires \texttt{m}, and
finishes at \texttt{z}.

\keepspace{11\baselineskip}
\begin{quizbox}[Can you construct an ordering that does not work?]
\textbf{Example broken schedule:}
\begin{center}
\sched{7/x, 8.5/y}{1/a, 2.5/b, 4/c}
\end{center}
The child runs first and signals at \texttt{b} while \emph{nobody is
waiting}, so the signal is lost. The parent then calls \texttt{wait()} at
\texttt{y} and \textbf{waits forever!}
\end{quizbox}

\subsection{Rule of thumb 1}

\begin{keybox}[Keep state in addition to CVs!]
\begin{itemize}
  \item CVs are used to \textbf{signal} threads when state \textbf{changes}.
  \item If the state is already as needed, the thread \textbf{doesn't wait}
        for a signal!
\end{itemize}
\end{keybox}

\keepspace{20\baselineskip}
\subsection{Attempt 2: add a \texttt{done} variable}

\begin{center}
\begin{minipage}[t]{0.48\textwidth}
\textbf{Parent:}
\begin{lstlisting}[numbers=none]
void thread_join() {
  Mutex_lock(&m);       // w
  if (done == 0)        // x
    Cond_wait(&c, &m);  // y
  Mutex_unlock(&m);     // z
}
\end{lstlisting}
\end{minipage}\hfill
\begin{minipage}[t]{0.48\textwidth}
\textbf{Child:}
\begin{lstlisting}[numbers=none]
void thread_exit() {
  done = 1;             // a
  Cond_signal(&c);      // b
}
\end{lstlisting}
\end{minipage}
\end{center}

\textbf{Fixes the previous broken ordering:}
\begin{center}
\sched{5/w, 6.5/x, 8/y, 9.5/z}{1/a, 2.5/b}
\end{center}
The child sets \texttt{done = 1} first. When the parent tests at
\texttt{x} it sees \texttt{done == 1}, so it does not wait at \texttt{y}
and goes straight to \texttt{z}.

\begin{quizbox}[Can you construct an ordering that does not work?]
\begin{center}
\sched{1/w, 2.5/x, 8/y}{4.5/a, 6/b}
\end{center}
\begin{enumerate}
  \item Parent: \texttt{w}, then \texttt{x} sees \texttt{done == 0}, so it
        is about to wait.
  \item Before it reaches \texttt{y}, the child runs: \texttt{a} sets
        \texttt{done = 1}, \texttt{b} signals --- but nobody is waiting yet.
  \item Parent executes \texttt{y}: \ldots\ \textbf{sleeps forever}\ldots
\end{enumerate}
The child does not take the mutex, so it can change \texttt{done} and
signal in the gap between the parent's \emph{check} and its \emph{wait}.
\end{quizbox}

\keepspace{24\baselineskip}
\subsection{Correct join}

\begin{center}
\begin{minipage}[t]{0.48\textwidth}
\textbf{Parent:}
\begin{lstlisting}[numbers=none]
void thread_join() {
  Mutex_lock(&m);       // w
  if (done == 0)        // x
    Cond_wait(&c, &m);  // y
  Mutex_unlock(&m);     // z
}
\end{lstlisting}
\end{minipage}\hfill
\begin{minipage}[t]{0.48\textwidth}
\textbf{Child:}
\begin{lstlisting}[numbers=none]
void thread_exit() {
  Mutex_lock(&m);       // a
  done = 1;             // b
  Cond_signal(&c);      // c
  Mutex_unlock(&m);     // d
}
\end{lstlisting}
\end{minipage}
\end{center}

\begin{center}
\sched{1/w, 2.5/x, 4/y, 10/z}{5.5/a, 7/b, 8.5/c}
\end{center}

\begin{keybox}[Use the mutex to ensure no race between interacting with state and wait/signal]
The child can only change \texttt{done} and signal while holding
\texttt{m}. The parent holds \texttt{m} from its check at \texttt{x} until
\texttt{wait()} releases it at \texttt{y}. So the child's \texttt{a--c}
cannot fall into the gap: either the child runs first (parent sees
\texttt{done == 1} and never waits), or the parent is already asleep when the
child signals (and is woken).
\end{keybox}

The slides use \texttt{if (done == 0)}. For this program, with one parent
and one child, \texttt{if} happens to be enough. The book still writes
\texttt{while (done == 0)} here: a \texttt{while} is \emph{always} correct,
an \texttt{if} only sometimes (Section~5 shows why).

% ===========================================================================
\section{The producer/consumer problem}
% ===========================================================================

\begin{itemize}
  \item \textbf{Producers} generate data (like pipe writers).
  \item \textbf{Consumers} grab data and process it (like pipe readers).
  \item Producer/consumer problems are frequent in systems, e.g.\ web
        servers.
\end{itemize}

\begin{keybox}[General strategy: use condition variables to]
\begin{itemize}
  \item make \textbf{producers wait} when buffers are \textbf{full};
  \item make \textbf{consumers wait} when there is \textbf{nothing to
        consume}.
\end{itemize}
\end{keybox}

\subsection{Bounded buffer}

A \textbf{bounded buffer} is used when you pipe the output of one program
into another.
\begin{itemize}
  \item Example: \texttt{grep foo file.txt | wc -l}
    \begin{itemize}
      \item the \texttt{grep} process is the \textbf{producer};
      \item the \texttt{wc} process is the \textbf{consumer};
      \item between them is an in-kernel \textbf{bounded buffer}.
    \end{itemize}
  \item The bounded buffer is a \textbf{shared resource} \(\Rightarrow\)
        \textbf{synchronized access} is required.
\end{itemize}

\keepspace{30\baselineskip}
\subsection{Example: UNIX pipes}

The slides animate a pipe's buffer. Shaded = data written but not yet read.
Data is read at \textbf{start} and written at \textbf{end}; both move to
the right and \textbf{wrap around} (the buffer is circular). Each panel is
the state after the action.

\begin{center}
\renewcommand{\arraystretch}{1.0}
\begin{tabular}{@{}r@{\hspace{0.6em}}l@{\hspace{1.5em}}r@{\hspace{0.6em}}l@{}}
  \textbf{0.} initial (empty) & \pipebuf{0/0}{0}{0} &
  \textbf{7.} read! (now empty) & \pipebuf{0/0}{0.93}{0.93} \\
  \textbf{1.} write! & \pipebuf{0/0.23}{0}{0.23} &
  \textbf{8.} read! & \pipebuf{0/0}{0.93}{0.93} \\
  \textbf{2.} write! & \pipebuf{0/0.6}{0}{0.6} &
  \textbf{9.} write! (wraps) & \pipebuf{0.93/1, 0/0.13}{0.93}{0.13} \\
  \textbf{3.} read! & \pipebuf{0.18/0.6}{0.18}{0.6} &
  \textbf{10.} write! & \pipebuf{0.93/1, 0/0.52}{0.93}{0.52} \\
  \textbf{4.} write! & \pipebuf{0.18/0.93}{0.18}{0.93} &
  \textbf{11.} write! (now full) & \pipebuf{0.93/1, 0/0.93}{0.93}{0.93} \\
  \textbf{5.} read! & \pipebuf{0.42/0.93}{0.42}{0.93} &
  \textbf{12.} write! & \pipebuf{0.93/1, 0/0.93}{0.93}{0.93} \\
  \textbf{6.} read! & \pipebuf{0.67/0.93}{0.67}{0.93} &
  \textbf{13.} read! & \pipebuf{0.07/0.93}{0.07}{0.93} \\
\end{tabular}
\end{center}

\begin{itemize}
  \item \textbf{Step 8} (buffer empty, \texttt{start = end}): a reader
        cannot proceed. \textbf{Note: readers must wait.}
  \item \textbf{Step 12} (buffer full, \texttt{end} has caught up with
        \texttt{start} from behind): a writer cannot proceed.
        \textbf{Note: writers must wait.}
  \item Step 13: a read frees space from \texttt{start} (which wraps past
        the right end), so writers can continue.
\end{itemize}

\subsection{Implementation requirements}

\begin{itemize}
  \item reads/writes to the buffer require \textbf{locking};
  \item when buffers are \textbf{full}, writers must \textbf{wait};
  \item when buffers are \textbf{empty}, readers must \textbf{wait}.
\end{itemize}

Start with the easy case:
\begin{itemize}
  \item 1 producer thread,
  \item 1 consumer thread,
  \item 1 shared buffer to fill/consume (max = 1).
\end{itemize}
\texttt{numfill} = number of buffers currently filled (\texttt{count} in the
code). We examine slightly broken code to begin\ldots

% ===========================================================================
\section{Single buffer: building the solution}
% ===========================================================================

\subsection{The put and get routines (version 1)}

\begin{lstlisting}
int buffer;
int count = 0; // initially, empty

void put(int value) {
     assert(count == 0);
     count = 1;
     buffer = value;
}

int get() {
     assert(count == 1);
     count = 0;
     return buffer;
}
\end{lstlisting}

\begin{itemize}
  \item Only put data into the buffer when \texttt{count} is zero, i.e.\
        when the buffer is \emph{empty}.
  \item Only get data from the buffer when \texttt{count} is one, i.e.\
        when the buffer is \emph{full}.
\end{itemize}

\keepspace{20\baselineskip}
\subsection{Producer/consumer threads (version 1)}

\begin{lstlisting}
void *producer(void *arg) {
     int i;
     int loops = (int) arg;
     for (i = 0; i < loops; i++) {
          put(i);
     }
}

void *consumer(void *arg) {
     int i;
     while (1) {
          int tmp = get();
          printf("%d\n", tmp);
     }
}
\end{lstlisting}

\begin{itemize}
  \item The producer puts an integer into the shared buffer \texttt{loops}
        number of times.
  \item The consumer gets the data out of that shared buffer.
\end{itemize}
There is no synchronization yet: the \texttt{assert}s fire as soon as the
producer calls \texttt{put()} on a full buffer or the consumer calls
\texttt{get()} on an empty one.

\keepspace{34\baselineskip}
\subsection{Single CV and \texttt{if} statement}

A single condition variable \texttt{cond} and associated lock
\texttt{mutex}.

\begin{lstlisting}
cond_t cond;
mutex_t mutex;

void *producer(void *arg) {
    int i;
    for (i = 0; i < loops; i++) {
        Pthread_mutex_lock(&mutex);           // p1
        if (count == 1)                       // p2
            Pthread_cond_wait(&cond, &mutex); // p3
        put(i);                               // p4
        Pthread_cond_signal(&cond);           // p5
        Pthread_mutex_unlock(&mutex);         // p6
    }
}

void *consumer(void *arg) {
    int i;
    for (i = 0; i < loops; i++) {
        Pthread_mutex_lock(&mutex);           // c1
        if (count == 0)                       // c2
            Pthread_cond_wait(&cond, &mutex); // c3
        int tmp = get();                      // c4
        Pthread_cond_signal(&cond);           // c5
        Pthread_mutex_unlock(&mutex);         // c6
        printf("%d\n", tmp);
    }
}
\end{lstlisting}

\begin{itemize}
  \item \texttt{p1--p3}: a producer waits for the buffer to be empty.
  \item \texttt{c1--c3}: a consumer waits for the buffer to be full.
  \item With just \emph{a single producer} and \emph{a single consumer},
        the code works.
\end{itemize}

\begin{quizbox}[If we have more than one of producer and consumer? What about 2 consumers?]
Can you find a problematic timeline with 2 consumers (still 1 producer)?
\end{quizbox}

\keepspace{26\baselineskip}
\subsection{Thread trace: broken solution (version 1)}

Consumers \Tc1, \Tc2 and producer \Tp; each row is one step, ``Count'' is
the buffer count after it.

\begin{center}
\small
\renewcommand{\arraystretch}{1.05}
\setlength{\tabcolsep}{4.5pt}
\begin{tabular}{@{}clclclcl@{}}
  \toprule
  \Tc1 & \textbf{State} & \Tc2 & \textbf{State} & \Tp & \textbf{State} & \textbf{Count} & \textbf{Comment} \\
  \midrule
  c1 & Running &    & Ready   &    & Ready   & 0 & \\
  c2 & Running &    & Ready   &    & Ready   & 0 & \\
  c3 & Sleep   &    & Ready   &    & Ready   & 0 & Nothing to get \\
     & Sleep   &    & Ready   & p1 & Running & 0 & \\
     & Sleep   &    & Ready   & p2 & Running & 0 & \\
     & Sleep   &    & Ready   & p4 & Running & 1 & Buffer now full \\
     & Ready   &    & Ready   & p5 & Running & 1 & \Tc1 awoken \\
     & Ready   &    & Ready   & p6 & Running & 1 & \\
     & Ready   &    & Ready   & p1 & Running & 1 & \\
     & Ready   &    & Ready   & p2 & Running & 1 & \\
     & Ready   &    & Ready   & p3 & Sleep   & 1 & Buffer full; sleep \\
     & Ready   & c1 & Running &    & Sleep   & 1 & \Tc2 sneaks in \ldots \\
     & Ready   & c2 & Running &    & Sleep   & 1 & \\
     & Ready   & c4 & Running &    & Sleep   & 0 & \ldots\ and grabs data \\
     & Ready   & c5 & Running &    & Ready   & 0 & \Tp awoken \\
     & Ready   & c6 & Running &    & Ready   & 0 & \\
  c4 & Running &    & Ready   &    & Ready   & 0 & \textbf{Oh oh! No data} \\
  \bottomrule
\end{tabular}
\end{center}

\begin{warnbox}[Why it breaks]
\begin{itemize}
  \item After the producer woke \Tc1, but \textbf{before \Tc1 ever ran}, the
        state of the bounded buffer was \textbf{changed by \Tc2}.
  \item There is \textbf{no guarantee} that when the woken thread runs, the
        state will still be as desired.
\end{itemize}
\Tc1 does not re-check \texttt{count} after waking (it used \texttt{if}), so
it calls \texttt{get()} on an empty buffer.
\end{warnbox}

\keepspace{30\baselineskip}
\subsection{Single CV and \texttt{while}}

Consumer \Tc1 wakes up and \textbf{re-checks} the state of the shared
variable. If the buffer is empty, the consumer simply goes back to sleep.
The only change is \texttt{if} \(\to\) \texttt{while} at \texttt{p2} and
\texttt{c2}:

\begin{lstlisting}
cond_t cond;
mutex_t mutex;

void *producer(void *arg) {
    int i;
    for (i = 0; i < loops; i++) {
        Pthread_mutex_lock(&mutex);           // p1
        while (count == 1)                    // p2
            Pthread_cond_wait(&cond, &mutex); // p3
        put(i);                               // p4
        Pthread_cond_signal(&cond);           // p5
        Pthread_mutex_unlock(&mutex);         // p6
    }
}

void *consumer(void *arg) {
    int i;
    for (i = 0; i < loops; i++) {
        Pthread_mutex_lock(&mutex);           // c1
        while (count == 0)                    // c2
            Pthread_cond_wait(&cond, &mutex); // c3
        int tmp = get();                      // c4
        Pthread_cond_signal(&cond);           // c5
        Pthread_mutex_unlock(&mutex);         // c6
        printf("%d\n", tmp);
    }
}
\end{lstlisting}

A simple rule to remember with condition variables is to \textbf{always use
while loops}.

\begin{keybox}[Why \texttt{while}: Mesa semantics]
Signaling a thread only \emph{wakes it up}. It is a \textbf{hint} that the
state has changed, but there is no guarantee that when the woken thread
actually runs, the state will still be as desired. This meaning of signal is
called \textbf{Mesa semantics}. The alternative, \textbf{Hoare semantics},
runs the woken thread immediately, which would make \texttt{if} safe, but it
is harder to build. Virtually every real system uses Mesa semantics, hence
\texttt{while}.
\end{keybox}

\subsection{Good rule of thumb 2}

\begin{keybox}[Whenever a lock is acquired, recheck assumptions about state!]
\begin{itemize}
  \item It is possible for another thread to grab the lock in between the
        signal and the wakeup from wait.
  \item Note that some libraries also have \textbf{``spurious wakeups''}
        (may wake multiple waiting threads at signal or at any time). Due
        to races inside the library's signal/wakeup code, two threads can
        wake up even though only one signal was sent. The \texttt{while}
        loop handles this too: the extra thread re-checks and sleeps again.
\end{itemize}
\end{keybox}

\begin{quizbox}[What about 2 consumers now?]
Can you find a problematic timeline with 2 consumers (still 1 producer)?
\end{quizbox}

\keepspace{36\baselineskip}
\subsection{Thread trace: single CV with \texttt{while} is still broken}

\begin{center}
\small
\renewcommand{\arraystretch}{1.05}
\setlength{\tabcolsep}{4.5pt}
\begin{tabular}{@{}clclclcl@{}}
  \toprule
  \Tc1 & \textbf{State} & \Tc2 & \textbf{State} & \Tp & \textbf{State} & \textbf{Count} & \textbf{Comment} \\
  \midrule
  c1 & Running &    & Ready   &    & Ready   & 0 & \\
  c2 & Running &    & Ready   &    & Ready   & 0 & \\
  c3 & Sleep   &    & Ready   &    & Ready   & 0 & Nothing to get \\
     & Sleep   & c1 & Running &    & Ready   & 0 & \\
     & Sleep   & c2 & Running &    & Ready   & 0 & \\
     & Sleep   & c3 & Sleep   &    & Ready   & 0 & Nothing to get \\
     & Sleep   &    & Sleep   & p1 & Running & 0 & \\
     & Sleep   &    & Sleep   & p2 & Running & 0 & \\
     & Sleep   &    & Sleep   & p4 & Running & 1 & Buffer now full \\
     & Ready   &    & Sleep   & p5 & Running & 1 & \Tc1 awoken \\
     & Ready   &    & Sleep   & p6 & Running & 1 & \\
     & Ready   &    & Sleep   & p1 & Running & 1 & \\
     & Ready   &    & Sleep   & p2 & Running & 1 & \\
     & Ready   &    & Sleep   & p3 & Sleep   & 1 & Must sleep (full) \\
  c2 & Running &    & Sleep   &    & Sleep   & 1 & Recheck condition \\
  c4 & Running &    & Sleep   &    & Sleep   & 0 & \Tc1 grabs data \\
  c5 & Running &    & Ready   &    & Sleep   & 0 & \textbf{Oops! Woke \Tc2} \\
  \midrule
  c6 & Running &    & Ready   &    & Sleep   & 0 & \textit{(cont.)} \\
  c1 & Running &    & Ready   &    & Sleep   & 0 & \\
  c2 & Running &    & Ready   &    & Sleep   & 0 & \\
  c3 & Sleep   &    & Ready   &    & Sleep   & 0 & Nothing to get \\
     & Sleep   & c2 & Running &    & Sleep   & 0 & \\
     & Sleep   & c3 & Sleep   &    & Sleep   & 0 & \textbf{Everyone asleep \ldots} \\
  \bottomrule
\end{tabular}
\end{center}

\begin{warnbox}[The problem: one CV for two kinds of waiters]
At \texttt{c5}, \Tc1 signals the single CV, which has both \Tc2 (a
consumer) and \Tp (a producer) waiting on it. It happens to wake \Tc2.
\Tc2 re-checks (thanks to \texttt{while}), sees an empty buffer and sleeps
again, and now nobody is left to wake the producer.

\textbf{A consumer should not wake other consumers, only producers, and
vice versa.}
\end{warnbox}

% ===========================================================================
\keepspace{24\baselineskip}
\section{The single-buffer producer/consumer solution}
% ===========================================================================

Use \textbf{two} condition variables and \texttt{while}:
\begin{itemize}
  \item \textbf{Producer} threads wait on the condition \texttt{empty}, and
        signal \texttt{fill}.
  \item \textbf{Consumer} threads wait on \texttt{fill} and signal
        \texttt{empty}.
\end{itemize}

\begin{lstlisting}
cond_t empty, fill;
mutex_t mutex;

void *producer(void *arg) {
    int i;
    for (i = 0; i < loops; i++) {
        Pthread_mutex_lock(&mutex);
        while (count == 1)
            Pthread_cond_wait(&empty, &mutex);
        put(i);
        Pthread_cond_signal(&fill);
        Pthread_mutex_unlock(&mutex);
    }
}
\end{lstlisting}

\keepspace{14\baselineskip}
\begin{lstlisting}[firstnumber=16]
void *consumer(void *arg) {
    int i;
    for (i = 0; i < loops; i++) {
        Pthread_mutex_lock(&mutex);
        while (count == 0)
            Pthread_cond_wait(&fill, &mutex);
        int tmp = get();
        Pthread_cond_signal(&empty);
        Pthread_mutex_unlock(&mutex);
        printf("%d\n", tmp);
    }
}
\end{lstlisting}

\begin{solbox}[Is this correct? Can you find a bad schedule?]
\textbf{Correct!}
\begin{itemize}
  \item no concurrent access to shared state (everything is under
        \texttt{mutex});
  \item every time the lock is acquired, assumptions are reevaluated
        (\texttt{while});
  \item a consumer will get to run after every \texttt{put()} (the producer
        signals \texttt{fill}, on which only consumers wait);
  \item a producer will get to run after every \texttt{get()} (the consumer
        signals \texttt{empty}, on which only producers wait).
\end{itemize}
\end{solbox}

\keepspace{10\baselineskip}
% ===========================================================================
\section{The final producer/consumer solution}
% ===========================================================================

More \textbf{concurrency} and \textbf{efficiency} \(\Rightarrow\) add more
buffer slots:
\begin{itemize}
  \item allow concurrent production or consuming to take place;
  \item reduce context switches.
\end{itemize}
Several values can now be produced before the producer has to sleep, and
several consumed before the consumer has to sleep. With just one producer
and one consumer, the gain is \textbf{fewer context switches}. With multiple
producers or consumers, they can also produce or consume
\textbf{concurrently}.

\keepspace{22\baselineskip}
\subsection{The final put and get routines}

\begin{lstlisting}
int buffer[MAX];
int fill = 0;
int use = 0;
int count = 0;

void put(int value) {
    buffer[fill] = value;
    fill = (fill + 1) % MAX;
    count++;
}

int get() {
    int tmp = buffer[use];
    use = (use + 1) % MAX;
    count--;
    return tmp;
}
\end{lstlisting}

This is the circular buffer of the UNIX-pipe example: \texttt{fill} plays
the role of \textbf{end} (where the next write goes), \texttt{use} the role
of \textbf{start} (where the next read comes from), both wrapping with
\texttt{\% MAX}. (The integer \texttt{fill} here is the buffer index; the
condition variable named \texttt{fill} below is a different variable.)

\keepspace{18\baselineskip}
\subsection{The final producer/consumer threads}

\begin{lstlisting}
cond_t empty, fill;
mutex_t mutex;

void *producer(void *arg) {
    int i;
    for (i = 0; i < loops; i++) {
        Pthread_mutex_lock(&mutex);
        while (count == MAX)
            Pthread_cond_wait(&empty, &mutex);
        put(i);
        Pthread_cond_signal(&fill);
        Pthread_mutex_unlock(&mutex);
    }
}
\end{lstlisting}

\keepspace{14\baselineskip}
\begin{lstlisting}[firstnumber=16]
void *consumer(void *arg) {
    int i;
    for (i = 0; i < loops; i++) {
        Pthread_mutex_lock(&mutex);
        while (count == 0)
            Pthread_cond_wait(&fill, &mutex);
        int tmp = get();
        Pthread_cond_signal(&empty);
        Pthread_mutex_unlock(&mutex);
        printf("%d\n", tmp);
    }
}
\end{lstlisting}

The only change from the single-buffer solution: the producer now waits
only when all \texttt{MAX} slots are full (\texttt{count == MAX}), so several
items can be produced before a consumer must run.

% ===========================================================================
\keepspace{32\baselineskip}
\section{Covering conditions}
% ===========================================================================

A memory allocator: threads ask for memory and wait until enough is free.

\begin{lstlisting}
// how many bytes of the heap are free?
int bytesLeft = MAX_HEAP_SIZE;

// need lock and condition too
cond_t c;
mutex_t m;

void *
allocate(int size) {
    Pthread_mutex_lock(&m);
    while (bytesLeft < size)
        Pthread_cond_wait(&c, &m);
    void *ptr = ...;          // get mem from heap
    bytesLeft -= size;
    Pthread_mutex_unlock(&m);
    return ptr;
}

void free(void *ptr, int size) {
    Pthread_mutex_lock(&m);
    bytesLeft += size;
    Pthread_cond_signal(&c); // whom to signal??
    Pthread_mutex_unlock(&m);
}
\end{lstlisting}

\keepspace{10\baselineskip}
\begin{quizbox}[Which waiting thread should be woken up?]
Assume there are zero bytes free.
\begin{itemize}
  \item Thread \(T_a\) calls \texttt{allocate(100)}.
  \item Thread \(T_b\) calls \texttt{allocate(10)}.
  \item Both \(T_a\) and \(T_b\) wait on the condition and go to sleep.
  \item Thread \(T_c\) calls \texttt{free(50)}.
\end{itemize}
\end{quizbox}

\begin{solbox}[Answer]
\(T_b\) should be woken: 50 free bytes satisfy its request for 10, but not
\(T_a\)'s request for 100. But \texttt{signal()} wakes \emph{some} waiting
thread, and the freeing thread cannot say which. If it wakes \(T_a\),
\(T_a\) re-checks, sees \(50 < 100\), and goes back to sleep, while
\(T_b\), who could have run, is never woken.
\end{solbox}

\keepspace{18\baselineskip}
\subsection{How to wake the right thread? Wake all the threads!}

\begin{keybox}[Solution (suggested by Lampson and Redell)]
\begin{itemize}
  \item Replace \texttt{pthread\_cond\_signal()} with
        \textbf{\texttt{pthread\_cond\_broadcast()}}.
  \item \texttt{pthread\_cond\_broadcast()}
    \begin{itemize}
      \item wakes up \textbf{all waiting threads};
      \item \textbf{cost}: too many threads might be woken;
      \item threads that shouldn't be awake will simply wake up, re-check
            the condition, and then go back to sleep.
    \end{itemize}
\end{itemize}
A condition variable used this way is called a \textbf{covering
condition}: the broadcast covers every case in which some thread needs to
wake up. Because every waiter re-checks its condition in a \texttt{while}
loop, waking extra threads is safe.
\end{keybox}

\begin{warnbox}[Broadcast vs.\ two CVs]
Broadcast would also have rescued the single-CV producer/consumer (the
woken consumer could no longer be the only one woken). There a better fix
existed: two CVs, so each side wakes only the other. As a rule, if your
program only works after changing signals to broadcasts (and you don't
think it should need to), you probably have a bug. In cases like the
allocator, where the signaler cannot know which waiter can proceed,
broadcast is the straightforward solution.
\end{warnbox}

% ===========================================================================
\section{Summary: rules of thumb for CVs}
% ===========================================================================

\begin{keybox}[Rules of thumb]
\begin{enumerate}
  \item \textbf{Keep state in addition to CVs.} CVs only signal changes; the
        state says whether there is anything to wait for (join's
        \texttt{done}, the buffer's \texttt{count}).
  \item \textbf{Always do wait/signal with the lock held.} Otherwise the
        state can change between the check and the wait (join attempt 2).
        For \texttt{wait} this is mandatory: \texttt{wait} assumes the lock
        is held, releases it while sleeping, and re-acquires it before
        returning. For \texttt{signal} it is not strictly needed in every
        program, but it is simplest and safest to always do it.
  \item \textbf{Whenever a thread wakes from waiting, recheck state}
        (\texttt{while}, not \texttt{if}): another thread may have run
        first, and there may be spurious wakeups.
\end{enumerate}
\end{keybox}

\begin{center}
\renewcommand{\arraystretch}{1.18}
\small
\begin{tabularx}{\textwidth}{@{}lX@{}}
  \toprule
  \textbf{Version} & \textbf{Problem / result} \\
  \midrule
  Join attempt 1 (no state) & Signal before wait is lost: parent waits forever. \\
  Join attempt 2 (state, child without lock) & Child can set \texttt{done} and signal between parent's check and wait. \\
  Join (correct) & State + lock held around check/wait and update/signal. \\
  P/C: single CV + \texttt{if} & Woken consumer finds buffer emptied by another consumer. \\
  P/C: single CV + \texttt{while} & Consumer can wake a consumer: everyone asleep. \\
  P/C: two CVs + \texttt{while} & Correct (single buffer). \\
  P/C: two CVs + \texttt{while}, \texttt{MAX} slots & Correct, more concurrency, fewer context switches. \\
  Allocator & Signal may wake the wrong waiter \(\Rightarrow\) use \texttt{broadcast} (covering condition). \\
  \bottomrule
\end{tabularx}
\end{center}

\end{document}
